\documentclass[]{article}
\usepackage{amsmath}
\usepackage{url}

\begin{document}
\begin{abstract}
Document to describe how the magnetic field is calculated
\end{abstract}

\section*{Summary of the spherical harmonic representions used to represent the magnetic field}
The REMIX horizontal current $\mathbf{J}(\theta, \phi)$ is represented as
\begin{equation}
\mathbf{J} = -\nabla \alpha + \mathbf{r}\times \nabla_S \Psi
\end{equation}
where $\nabla_S$ is the surface gradient operator. $\Psi$ and $\alpha$ are expanded in global spherical harmonics. We follow the convention for spherical harmonic coefficients in [Laundal et al. 2025] and write
\begin{align}
\Psi   = & \sum_{n = 1}^{N}\sum_{m = 0}^{\mathrm{min}(n, M)} P_n^m(\cos\theta)[\kappa_n^m'\cos(m\phi) + l_n^m'\sin(m\phi)] \nonumber \\
=& -\frac{1}{\mu_0}\sum_{n = 1}^{N}\sum_{m = 0}^{\mathrm{min}(n, M)} P_n^m(\cos\theta)(2n+1)[\kappa_n^m\cos(m\phi) + l_n^m \sin(m\phi)] \\
\alpha = & \sum_{n = 1}^{N}\sum_{m = 0}^{\mathrm{min}(n, M)} P_n^m(\cos\theta)[\psi_n^m'\cos(m\phi) + \xi_n^m' \sin(m\phi)] \nonumber \\
=& -\frac{1}{\mu_0}\sum_{n = 1}^{N}\sum_{m = 0}^{\mathrm{min}(n, M)} P_n^m(\cos\theta)n(n+1)[\psi_n^m\cos(m\phi) + \xi_n^m \sin(m\phi)].
\end{align}
where we choose $R$ to represent the radius of the ionosphere\footnote{note that this is different from the convention in geomagnetism that $R=6371.2$~km.}. The script \path{sh_analysis_lompeosse_inverse.py} calculates the expansion coefficients $\kappa_n^m'$, $l_n^m'$, $\psi_n^m'$, and $\xi_n^m'$. From the equations above, we find that
\begin{align}
\{\kappa_n^m, l_n^m\} =& -\frac{\mu_0}{2n + 1}\{\kappa_n^m', l_n^m'\} \\
\{\psi_n^m, \xi_n^m\} =& -\frac{\mu_0}{n(n+1)}\{\psi_n^m', \xi_n^m'\}.
\end{align}
These coefficients will be used to calculate the magnetic field on ground and in space.

\subsection*{Ground magnetic field}
On ground we can assume\footnote{if the main magnetic field is radial, which is an ok approximation at high latitudes} that we only see the part of the current represented by $\Psi$. We have that
\begin{equation}
\mathbf{B}(r < R, \theta, \phi) = -\nabla \sum_{n,m}(n+1)\left(\frac{r}{R}\right)^n P_n^m(\theta)[\kappa_n^m\cos(m\phi) + l_n^m \sin(m\phi)]
\end{equation}
or
\begin{equation}
\mathbf{B}(r < R, \theta, \phi) = \mu_0\nabla \sum_{n,m}\frac{n+1}{2n+1}\left(\frac{r}{R}\right)^n P_n^m(\theta)[\kappa_n^m'\cos(m\phi) + l_n^m' \sin(m\phi)]
\end{equation}

\subsection*{Space magnetic field}
In space we see both the part of the current represented by $\Psi$, $\mathbf{B}_\Psi$, (the divergence-free part of the horizontal current) and the part of the current represented by $\alpha$, $\mathbf{B}_\alpha$, (the curl-free part of the horizontal current + its source, the field-aligned current). We have that the first part is 
\begin{equation}
\mathbf{B}_\Psi(r > R, \theta, \phi) = \nabla \sum_{n,m}n\left(\frac{R}{r}\right)^{n+1} P_n^m(\theta)[\kappa_n^m\cos(m\phi) + l_n^m \sin(m\phi)]
\end{equation}
or
\begin{equation}
\mathbf{B}_\Psi(r > R, \theta, \phi) = -\mu_0\nabla \sum_{n,m}\frac{n}{2n+1}\left(\frac{R}{r}\right)^{n+1} P_n^m(\theta)[\kappa_n^m'\cos(m\phi) + l_n^m' \sin(m\phi)].
\end{equation}

The $\alpha$-part of the magnetic field can be written as
\begin{equation}
\mathbf{B}_\alpha(r > R, \theta, \phi) = \hat{\mathbf{r}} \times \nabla \sum_{n,m}P_n^m(\theta)[\psi_n^m\cos(m\phi) + \xi_n^m\sin(m\phi)]
\end{equation}
or
\begin{equation}
\mathbf{B}_\alpha(r > R, \theta, \phi) = -\hat{\mathbf{r}} \times \nabla \sum_{n,m}\frac{\mu_0}{n(n+1)}P_n^m(\theta)[\psi_n^m'\cos(m\phi) + \xi_n^m'\sin(m\phi)]
\end{equation}

\subsection*{Some remarks}
The magnetic field resolution will be restricted by the resolution of the spherical harmonics. This means that the OSSE should only be used for scale sizes that are larger than $\sim 2\pi r/(N + 1/2)$, which is unfortunately quite large-scale. For example, if we go to $N=200$ (which is high!) we can resolve $\sim 200$~km structures, while $N=100$ gives $\sim 400$~km spatial resolution. 

Another approximation that is made here is that the field lines are radial. We could include the effect of inclined magnetic field lines by using the same approach as in [Laundal et al. 2025], but then we would need to also save several quite large matrices (hundreds of MB) to map from the $\alpha$ coefficients to coefficients that describe potential magnetic fields. It would be interesting to include this at a later point.

We ignore ground induction in these calculations. Ground induction effects could be approximated by for example the mirror current method. This would change the $\Psi$ magnetic field. 



\end{document}